% ============================================
% DEMO SCENARIOS & WALKTHROUGH
% ============================================
% 👤 PERSON 4: Implementation & Demo Experience
% 
% Instructions:
% - Describe 2-3 concrete demo scenarios attendees will experience
% - Walk through typical interaction flow
% - Explain what makes the demo compelling/educational
% - Reference demo_complete.py functionality
% - THIS IS CRITICAL FOR DEMO TRACK - make it vivid
% - Length: 0.5-0.7 pages
% ============================================

\section{Demo Scenarios and Walkthrough}

[PERSON 4: Paint a picture of what attendees will see and do]

\subsection{Interactive Demo Experience}

During the live demonstration, attendees engage with ADRDS through hands-on 
exploration of ransomware detection scenarios. The \texttt{demo\_complete.py} 
script orchestrates the following interactive experiences:

\subsection{Scenario 1: Threshold Exploration}

Attendees adjust ML classification thresholds in real-time to observe the trade-off 
between detection rate and false positive rate. As ransomware behaviors progress 
through stages, users see:
\begin{itemize}
    \item Detection confidence scores updating dynamically
    \item Visual alerts triggering when thresholds are exceeded
    \item Impact on system performance metrics
\end{itemize}

\textbf{Learning Outcome:} Understanding the practical implications of threshold 
tuning in production security systems.

\subsection{Scenario 2: Attack Stage Progression}

Users control the pace and characteristics of simulated ransomware attacks:
\begin{itemize}
    \item Select target file types and encryption patterns
    \item Enable/disable C2 communication
    \item Trigger persistence mechanisms
\end{itemize}

The system visualizes behavioral indicators at each stage, showing which features 
contribute most to detection decisions.

\textbf{Learning Outcome:} Recognizing stage-specific behavioral signatures and 
understanding multi-stage detection strategies.

\subsection{Scenario 3: Model Comparison}

Attendees compare performance of different ML classifiers (Random Forest, SVM, 
Neural Network) on the same behavioral sequence:
\begin{itemize}
    \item Side-by-side detection timelines
    \item Comparative accuracy and latency metrics
    \item Feature importance rankings per model
\end{itemize}

\textbf{Learning Outcome:} Evaluating ML model trade-offs for real-world deployment.

\subsection{Demo Flow}

A typical 10-minute demo session follows this structure:
\begin{enumerate}
    \item Brief system overview (2 min)
    \item Scenario 1 walkthrough with attendee participation (3 min)
    \item Scenario 2 exploration (3 min)
    \item Q\&A and custom scenario testing (2 min)
\end{enumerate}

\textbf{Interactivity:} Attendees can propose custom attack scenarios, adjust 
any simulation parameter, and observe immediate detection outcomes, making each 
demo session unique and engaging.
